\documentclass[11pt]{article}
\input{../_estilo.tex}
\renewcommand{\tituloejercicio}{Tema 1 --- Ejercicio 9}

\begin{document}

{\LARGE\bfseries\color{azultitulo} Ejercicio 9 --- Una mezcla ``optimizada'' que deja de comparar}

\vspace{2pt}
\textbf{Dificultad:} $\bigstar\bigstar$ \quad \textbf{Conferencias:} C5 (Cotas mínimas)
\vspace{8pt}
\hrule
\vspace{12pt}

\section*{Enunciado}
\begin{enunciado}
Una mezcla ``optimizada'' de dos listas ordenadas de largo $n$ deja de comparar cuando a la primera lista
le queda un solo elemento, y lo pone al final. Cuenta cuántas comparaciones hace en el peor caso. Sin
probarla, usa la sección 6 de la conferencia 5 para decidir si es correcta, y construye con el adversario
la entrada en la que falla.
\end{enunciado}

\section{El algoritmo: una mezcla con parada prematura}

La variante intenta ahorrar una comparación: mezcla normalmente mientras a la primera lista le queden al
menos dos elementos, y cuando queda uno supone que puede colocarlo al final. La pregunta importante no es
solo cuánto trabajo ahorra, sino si esa suposición conserva la especificación de una mezcla ordenada.

\textbf{Especificación.} Precondición: $A$ y $B$ están ordenados de forma creciente y tienen el mismo largo
$n\ge1$. Postcondición: devuelve una lista ordenada que contiene exactamente una copia de cada elemento de
$A$ y de $B$.

\lstinputlisting[style=python,firstline=14,lastline=28]{verif.py}

\section{Correctitud}

Como en los ejercicios 7 y 8, separamos correctitud parcial, terminación y correctitud total. Aquí el
resultado será distinto: la terminación sí se puede demostrar, pero la correctitud parcial es falsa. La
cota mínima de C5 permite detectarlo sin comenzar por una batería de pruebas.

\subsection*{Correctitud parcial: qué tendría que garantizar}

Para que la postcondición se cumpla, no basta con que $R$ contenga todos los elementos: también debe quedar
ordenada. El código sí conserva el multiconjunto de entrada, porque cada elemento se copia una sola vez,
pero cuando el ciclo termina por dejar un elemento de $A$, coloca todo $B[j:]$ antes de $A[i:]$ sin comparar
sus fronteras. Por tanto, queda por demostrar si ese orden es siempre válido; la siguiente cota muestra que
no puede serlo.

\subsection*{La cota mínima como detector de errores}

\begin{resultado}{C5 §6, el teorema que ya tenemos}
Mezclar dos listas ordenadas de largo $n$ requiere, en el peor caso, \textbf{exactamente} $2n-1$
comparaciones --- ningún algoritmo correcto hace menos.
\end{resultado}

La variante optimizada hace a lo sumo $2n{-}2<2n{-}1$ comparaciones en \textbf{cualquier} entrada (se demuestra
en la sección~\ref{sec:costo}). Si fuera
correcta \textbf{para toda entrada}, contradiría directamente el teorema. Por lo tanto, \textbf{tiene que
haber alguna entrada en la que se equivoca}: lo sabemos sin ejecutar el algoritmo, comparando únicamente
$2n{-}2$ contra $2n{-}1$.

Este es exactamente el uso de una cota mínima como \textbf{detector de errores} que describe C5 §8: una cota
inferior no solo mide rendimiento, también restringe lo que puede hacer cualquier algoritmo correcto. Si un
algoritmo de comparaciones resuelve siempre el problema pero usa menos comparaciones que la cota en todas las
entradas, hay una contradicción.

\subsection*{Terminación}

Sean $i$ y $j$ las cantidades de elementos ya consumidos de $A$ y $B$. Mientras el ciclo está activo se
cumple $i\le n-2$ y $j\le n-1$. Usamos el potencial
\[
\Phi=(n-1-i)+(n-j)=2n-1-(i+j).
\]
Es un entero no negativo; cada vuelta consume exactamente un elemento, de $A$ o de $B$, y por tanto reduce
$\Phi$ en 1. El ciclo termina cuando a $A$ le queda un solo elemento o cuando $B$ se vacía. Las operaciones
\texttt{extend} posteriores recorren listas finitas, así que el procedimiento completo termina.

\subsection*{Correctitud total}

La terminación está garantizada, pero la postcondición no: la construcción adversaria de la sección~\ref{sec:adv}
produce una salida que contiene todos los elementos, aunque no está ordenada. Por eso la versión optimizada
no tiene correctitud total. La cota mínima demuestra que el fallo debe existir; el adversario muestra
exactamente dónde aparece.

\section{Costo: cuántas comparaciones hace}\label{sec:costo}

Cada vuelta del ciclo hace una comparación y consume un elemento. Si el ciclo termina porque a $A$ le queda
uno, se han consumido como máximo $n-1$ elementos de $A$ y, como $j<n$, como máximo $n-1$ de $B$. Si termina
porque $B$ se vacía, entonces se han consumido $n$ elementos de $B$ y, antes de la última comparación,
quedaban al menos dos elementos de $A$, de modo que se consumieron como máximo $n-2$ de $A$. En ambos casos,
el número de comparaciones es como máximo
\[
(n-1)+(n-1)=2n-2.
\]

La mezcla correcta también tarda $O(n)$ tiempo, pero su peor caso necesita $2n-1$ comparaciones. La variante
ahorra una comparación como máximo, a cambio de perder la garantía de orden. La salida ocupa $O(n)$ espacio;
si no contamos el arreglo de salida, el algoritmo con índices necesita espacio auxiliar $O(1)$ (el listado de
Python copia la cola con las rebanadas \texttt{B[j:]} y \texttt{A[i:]}, que añade $O(n)$ temporal).

\textbf{Verificado} con un barrido aleatorio de 300 instancias por tamaño:
\begin{lstlisting}[style=salida]
     n  max comp. optimizada   2n-2  max comp. correcta   2n-1
     2                     2      2                   3      3
     5                     8      8                   9      9
    10                    18     18                  19     19
    50                    98     98                  99     99
   200                   398    398                 399    399
  1000                  1998   1998                1999   1999
\end{lstlisting}
En cada tamaño medido la variante alcanza $2n-2$ (la columna ``max comp. optimizada'' coincide con la de $2n-2$), así que la
cota no es floja, y la mezcla correcta, medida con el mismo barrido, alcanza $2n-1$, como dice C5 §6.

\section{Construcción del adversario: la entrada en la que falla}\label{sec:adv}

C5 §6 construye la cota $2n{-}1$ con la entrada \textbf{intercalada perfecta}
$a_1<b_1<a_2<b_2<\cdots<a_n<b_n$ --- la reutilizamos íntegra, sin cambiar nada, porque es exactamente la
entrada diseñada para que un algoritmo de mezcla tenga que comparar \emph{todos} los pares de vecinos.

\begin{lstlisting}[style=salida]
  n=4: A=[1, 3, 5, 7] B=[2, 4, 6, 8]
       salida optimizada: [1, 2, 3, 4, 5, 6, 8, 7]
       correcta (ordenada)?: False  (deberia ser: False)
  n=5: A=[1, 3, 5, 7, 9] B=[2, 4, 6, 8, 10]
       salida optimizada: [1, 2, 3, 4, 5, 6, 7, 8, 10, 9]
       correcta (ordenada)?: False  (deberia ser: False)
  n=8: A=[1, 3, 5, 7, 9, 11, 13, 15] B=[2, 4, 6, 8, 10, 12, 14, 16]
       salida optimizada: [1, 2, 3, 4, 5, 6, 7, 8, 9, 10, 11, 12, 13, 14, 16, 15]
       correcta (ordenada)?: False  (deberia ser: False)
\end{lstlisting}

\textbf{Por qué falla, exactamente.} Con esta entrada, la mezcla normal va alternando $a_1,b_1,a_2,b_2,
\dots$ Justo cuando a $A$ le queda un solo elemento ($a_n$), ya se consumieron $a_1,\dots,a_{n-1}$ y
$b_1,\dots,b_{n-1}$ --- a $B$ le queda un solo elemento, $b_n$. El atajo pone ese elemento de
$B$ primero y después, al final de todo, el $a_n$ que quedaba pendiente. Pero en la cadena,
$b_{n-1}<a_n<b_n$: el orden correcto termina en $\dots,b_{n-1},a_n,b_n$, y el atajo produce
$\dots,b_{n-1},b_n,a_n$ --- $a_n$ y $b_n$ quedan invertidos. Se ve exactamente así en la salida de $n=4$:
termina en $\dots,6,8,7$ en vez de $\dots,6,7,8$.

\section{Generadores dirigidos y datos extra}

El contraejemplo intercalado demuestra el fallo, pero una sola entrada no muestra qué ocurre en otras
formas de detener el ciclo. Por eso \texttt{verif.py} incluye generadores dirigidos: algunos son controles
en los que la parada prematura sí es segura y otros producen distintos tipos de salida rota. Además de
comprobar si la salida está ordenada, se cuentan sus \textbf{inversiones}, es decir, pares de elementos que
aparecen en el orden contrario al correcto. Así medimos no solo si falla, sino cuánto se desordena la salida.

\begin{center}
\scriptsize
\renewcommand{\arraystretch}{1.16}
\begin{tabular}{>{\raggedright\arraybackslash}p{.20\linewidth}
                >{\raggedright\arraybackslash}p{.30\linewidth}
                >{\raggedright\arraybackslash}p{.32\linewidth}}
\toprule
\textbf{Generador} & \textbf{Construcción} & \textbf{Qué comprueba} \\
\midrule
\texttt{gen\_B\_se\_vacia} & B contiene valores menores que todos los de A. & B se consume completamente; la parada ocurre por la condición normal y la salida es correcta. \\
\texttt{gen\_parada\_correcta} & Quedan $r$ valores de B, todos menores que el último de A. & La variante deja de comparar, pero en este caso concreto colocar B antes del último de A sí es seguro. \\
\texttt{gen\_una\_inversion} & Queda un valor de B mayor que el último de A. & Produce una única inversión al copiar ese valor antes del último de A. \\
\texttt{gen\_varias\_inversiones} & Quedan $r$ valores de B, todos mayores que el último de A. & Produce $r$ inversiones; muestra que el error puede crecer con el tamaño de la cola. \\
\texttt{gen\_cola\_mixta} & La cola de B contiene valores menores y mayores que el último de A. & Separa los valores que deberían ir antes y después, pero el atajo los copia todos antes. \\
\texttt{gen\_intercalada} & $A=(1,3,\ldots,2n-1)$ y $B=(2,4,\ldots,2n)$. & Es el adversario de C5 y fuerza la inversión final clásica. \\
\texttt{gen\_aleatoria\_ordenada} & Dos listas ordenadas obtenidas con semillas fijas. & Mide la frecuencia del fallo en entradas no diseñadas a mano. \\
\bottomrule
\end{tabular}
\end{center}

\clearpage
La salida resumida de la batería es la siguiente. ``Comp. media'' y ``comp. máx.'' cuentan las
comparaciones; ``inv. media'' y ``inv. máx.'' cuentan las inversiones de la salida:

\begin{lstlisting}[style=salida]
               generador  casos   fallos  comp.media  comp.max  inv.media  inv.max
              B se vacia     29    0.00%       16.00        30       0.00        0
         parada correcta     28    0.00%       29.04        56       0.00        0
           una inversion     29  100.00%       29.00        57       1.00        1
      varias inversiones     25  100.00%       30.00        54       5.00        5
              cola mixta     26  100.00%       30.00        55       2.00        2
  intercalada adversaria     29  100.00%       29.00        57       1.00        1
    aleatorias ordenadas   1000   51.60%       28.63        58       0.91        7
\end{lstlisting}

Estos resultados añaden tres matices a la demostración. Primero, no toda parada temprana falla: los dos
generadores de control tienen 0\% de fallos. Segundo, cuando falla, la salida puede tener una inversión, varias, o solo
las de los valores de la cola que superan al último de $A$ (cola mixta), según la cola que quedó sin comparar. Tercero, en 1000 entradas
aleatorias ordenadas el fallo aparece en el 51.60\% de los casos; por tanto, el adversario no es un fenómeno
aislado, aunque la prueba de la cota mínima sigue siendo la demostración general.

\begin{lstlisting}[style=salida]
ejemplos de salidas rotas:
  una inversion: A=[0, 2, 4, 6, 8, 10, 12, 16], B=[1, 3, 5, 7, 9, 11, 13, 17]
    salida=[..., 10, 11, 12, 13, 17, 16], comparaciones=13, inversiones=1
  varias inversiones: A=[0, 2, 4, 6, 8, 10, 12, 16], B=[1, 3, 5, 7, 9, 17, 18, 19]
    salida=[..., 10, 12, 17, 18, 19, 16], comparaciones=12, inversiones=3
  cola mixta: A=[0, 2, 4, 6, 8, 10, 12, 19], B=[1, 3, 5, 7, 17, 18, 21, 22]
    salida=[..., 12, 17, 18, 21, 22, 19], comparaciones=11, inversiones=2
  intercalada adversaria: A=[1, 3, 5, 7], B=[2, 4, 6, 8]
    salida=[..., 3, 4, 5, 6, 8, 7], comparaciones=5, inversiones=1
\end{lstlisting}

\section*{Respuesta a la pregunta del enunciado}

\begin{resultado}{Conclusión}
\begin{itemize}[leftmargin=14pt]
\item \textbf{Comparaciones:} la variante hace como máximo $2n-2$ comparaciones, una menos que la cota
mínima $2n-1$ de una mezcla correcta en el peor caso.
\item \textbf{Correctitud:} no es correcta para todas las entradas. La cota de C5 permite saberlo sin
probarla, porque un algoritmo correcto no puede resolver siempre la mezcla usando menos que la cota mínima.
\item \textbf{Adversario:} las listas intercaladas $A=[1,3,\ldots,2n-1]$ y
$B=[2,4,\ldots,2n]$ hacen que el algoritmo termine con $\dots,2n-2,2n,2n-1$ en lugar de
$\dots,2n-2,2n-1,2n$.
\item \textbf{Terminación y costo:} el procedimiento sí termina, tarda $O(n)$ y usa $O(1)$ espacio
auxiliar con índices, aparte de la lista de salida, pero el ahorro de una comparación destruye la postcondición.
\end{itemize}
\end{resultado}

\section*{Errores que hay que vigilar}

\begin{alerta}{}
\begin{itemize}[leftmargin=14pt]
\item \textbf{Pensar que el error es raro, como los de las conferencias 2 y 5.} Los errores plantados de
C2 §7 (cerca del 1\% de las instancias pequeñas y ninguna de las grandes) y de C5 §8 (probabilidad $1/(n-1)$,
0.12\% con $n=1024$) casi no se ven con entradas aleatorias --- hace falta un adversario cuidadoso para
encontrarlos. \textbf{Este no es de esa clase}: medido sobre 20{,}000 pares de arreglos aleatorios
uniformes (tamaños entre 2 y 30), falla en
\begin{lstlisting}[style=salida]
entradas aleatorias uniformes: 9834/20000 fallan (49.17%)
\end{lstlisting}
casi la \textbf{mitad} de las veces. La razón es estructural, no un caso extremo: el atajo falla cada vez
que, en el momento en que $A$ queda con un solo elemento, a $B$ todavía le queda algún elemento \emph{mayor}
que ese último de $A$ (equivale a $b_n>a_n$, es decir, a que el máximo de $B$ sea mayor que el de $A$) --- y con datos
genéricos eso pasa casi la mitad de las veces, no es lo raro. La lección de C2 §6 (el 7 ``es una garantía y
no una descripción'': el caso típico no delata el error) también vale al revés: a veces el caso típico
\emph{sí} delata el error, y una batería de pruebas al azar basta para sospecharlo.
\item \textbf{Confundir ``hace menos comparaciones'' con ``es más rápido y sigue siendo correcto''.}
Reducir el conteo de comparaciones, en el peor caso, por debajo de una cota demostrada para el problema
\textbf{siempre} significa que el algoritmo ya no resuelve el problema en todas las entradas (o que la cota no aplica
porque cambió el modelo --- no es el caso aquí, la mezcla sigue siendo un algoritmo de comparaciones puro).
\item \textbf{Arreglarlo mal, comparando una vez más ``por si acaso''.} Agregar una comparación extra entre
el último elemento de $A$ y el último de $B$ antes de hacer \texttt{extend} no alcanza en general: si a $B$
le quedan \textbf{varios} elementos mayores que el último de $A$ (no solo el último de $B$), una sola
comparación adicional no alcanza para ubicarlo entre ellos. El arreglo correcto es simplemente no cortar el ciclo principal hasta que uno de
los dos arreglos se vacíe del todo --- la mezcla estándar de C5 §6, sin atajos.
\end{itemize}
\end{alerta}

\end{document}
